defining @@M0@@, and let @@M1@@ be the generic point of @@M2@@. Since @@M3@@ is integral, @@M4@@ is the maximal ideal of @@M5@@, and since @@M6@@ is regular, @@M7@@ can be generated by an @@M8@@-regular sequence. We can therefore find a sequence of sections of @@M9@@ in a neighborhood of @@M10@@ forming an @@M11@@-regular sequence. Consequently, there is an open neighborhood @@M12@@ of @@M13@@ such that @@M14@@ is contained in @@M15@@, and @@M16@@ is a regular ideal of @@M17@@. The immersion @@M18@@ is then a regular immersion, with codimension @@M19@@ (EGA IV 5.1.2). Furthermore, we have the exact sequence (IV\qquad):
@@M20@@
where @@M21@@ is the immersion @@M22@@ and @@M23@@ the immersion @@M24@@; tensoring with @@M25@@ gives an analogous sequence. Now @@M26@@, so by 4.6 the image @@M27@@ of the class of @@M28@@ under @@M29@@ belongs to @@M30@@. Since @@M31@@ is a surjective @@M32@@-homomorphism, it induces a \emph{surjection} @@M33@@, so @@M34@@ with @@M35@@; there is therefore @@M36@@ such that:
@@M37@@
Now @@M38@@, and since @@M39@@, @@M40@@. Consequently, @@M41@@ can be expressed as a combination of classes of sheaves @@M42@@, where the @@M43@@ are integral closed subschemes of @@M44@@, distinct from @@M45@@. By the induction hypothesis, @@M46@@, which completes the proof.

\sgaStableTarget{sga6.e07.bibliography}\subsection*{BIBLIOGRAPHY}
\begin{enumerate}[label={[\arabic*]}]
\item Cartan--Eilenberg, Homological Algebra.
\end{enumerate}



% Expose VIII macro additions
\providecommand{\etale}{\'etale}
\providecommand{\Ch}{\operatorname{Ch}}


